\subsection{Explanatory target and vocabulary follow-up}\label{sec:followup}
The final adaptive follow-up contains twelve new Swin-T runs: a seven-output minority-permutation target at seeds 17, 42 and 89, plus eight-output Soft, Uniform and Tie-Uniform at the same seeds. Training conditions, checkpoint rules and the primary analysis plan were frozen before full training, after the earlier manuscript and review findings were known. A separate ranking extension was frozen while training was underway, before any new calibration or test inference; original outcomes were already known. Every new epoch's weights and selection predictions are retained locally; all twelve runs finished and were hashed before new calibration or test inference. No condition was added or omitted according to its result.

The permutation target fixes every category position attaining the maximum. Remaining positions are cycled in a row-specific shuffled order, using PCG64 seeded from the SHA-256 digest of \texttt{minority-permutation-20261020:} and the original row ID. This mapping stays fixed across epochs and training seeds. It preserves each row's probability multiset, entropy and all maximum positions, while disrupting minority-category assignment when unequal values move. Rows with fewer than two nonmaximum positions are unchanged; duplicate probabilities can also yield an unchanged target. It changes 20,308 of 28,709 Training targets (70.74\%; 71.11\% of positive-mass rows). Its weighted class-marginal L1 difference from Soft is 0.1242 (Fig.~\ref{fig:targets}). Aggregate priors are therefore not preserved, limiting causal interpretation despite matching per-image entropy. Seven-output initialization is checked against the retained original runs.

The expanded vocabulary appends contempt as the eighth expression output; unknown and not-a-face remain outside the classifier and remain in the ten-category reference. Supported mass and targets are recomputed for $K=8$, leaving 150 zero-mass Training rows. Soft retains $r$, Uniform assigns $(1-m)/(K-1)$ off its chosen maximum, and Tie-Uniform assigns $(1-km)/(K-k)$ off its $k$ tied maxima when $k<K$. When all $K$ positions tie, Tie-Uniform assigns $1/K$ to each. Within each eight-output seed, targets share initialization, sample order, augmentation and the original twelve-epoch budget. The vocabulary comparison changes output dimension and supported weights together; it evaluates that policy change rather than a uniquely isolated contempt effect.

Distance-selected and epoch-12 states are reported for all new runs. A predeclared additional rule selects the earliest epoch minimizing complete-vote disagreement at 80\% MSP coverage on the selection subset for every eight-output condition. It tests a different stopping objective without choosing on test outcomes. This alternative cannot be reconstructed for all historical seven-output runs, whose intermediate weights were not retained.

The principal descriptive follow-up is distance-selected eight-output Soft minus Uniform at 80\% own-MSP coverage, with all three paired differences, sample standard deviation and conditional pixel-cluster intervals. The fixed epoch-12 comparison is a planned budget-matched sensitivity; the separately risk-selected comparison is an additional stopping sensitivity, not a replacement chosen for its test result. Tie-Uniform and permutation comparisons, full-coverage distance, fixed-set cross-evaluation, risk components and coverage curves are retained regardless of direction. These analyses extend the same exposed source rather than repair prior test reuse. Reported intervals condition on the fitted models and annotations and do not establish generalization over future training seeds.
